\documentclass{article}
\usepackage[utf8]{inputenc}
\usepackage{amsmath, amssymb}
\usepackage[a4paper, top=2.5cm, bottom=2.5cm, left=3cm, right=3cm]{geometry}

\title{Number and Group Theory\\\textbf{TC50331E}\\Algebraic Coding Theory}
\author{Jaish Singh\\33139652}
\date{April 2026}

\begin{document}
\maketitle
\setlength{\parskip}{0.5cm}

\section{Block Codes}

%-----------------------------  1 (a)  -----------------------------
Let \( \mathbf{B}(n)=\mathbb{Z}_2^n \), the set of all binary strings of length n. That is,
\[
    \mathbf{B}(n) = \{ (x_1, \dots, x_n) \mid x_i \in \mathbb{Z}_2 \text { for }
  1 \le i \le n \}.
\]
with a binary component-wise addition modulo 2: 
\( (x_1, \dots, x_n) \oplus (y_1, \dots, y_n) = (x_1 + y_1, \dots, x_n + y_n), \)
where addition is taken in $\mathbb{Z}_2$. This operation is well-defined 
since \( \mathbb{Z}_2 \) is closed under addition. We verify the group axioms 
for \((\mathbf{B}(n), \oplus )\):
\begin{enumerate}
    \item \textit{Closure: } Let \( x, y \in \mathbf{B}(n) \). Then \( x \oplus y \in \mathbf{B}(n) \) 
        since addition in \( \mathbb{Z}_2 \) is closed.

    \item \textit{Associativity: } For all \( x, y, z \in \mathbf{B}(n) \), we have
    \[ (x \oplus y) \oplus z = x \oplus (y \oplus z), \] since addition in
    \( \mathbb{Z}_2 \) is associative.

\item \textit{Identity: } The element \( 0 = (0, \dots, 0) \in \mathbf{B}(n) \) satisfies 
    \( x \oplus 0 = x \) for all \( x \in \mathbf{B}(n) \).

\item \textit{Inverses: } For each \( x \in \mathbf{B}(n) \), we have \( x \oplus x = 0 \),
        so every element is its own inverse.
\end{enumerate}
Having established that \( (\mathbf{B}(n), \oplus) \) is a group, we now observe that it is in
fact abelian. This follows from the commutativity of addition in \( \mathbb{Z}_2 \), 
applied component-wise. Since digital signals are inherently binary, it is natural to 
model transmitted data as elements of \(\mathbb{Z}_2 = \{0, 1\}\), the group of integers
modulo 2 under addition. The group operation, which coincides with XOR, captures parity 
computations precisely: the parity of a binary string is simply its sum in \(\mathbb{Z}_2\).
This structure extends naturally to \(\mathbb{Z}_2^n\), the group of binary strings of 
length n, which serves as the space in which code words live. 

%-----------------------------  1 (b)  -----------------------------
For \(0 < k < n\), we define an (n, k) \textit{block code} to be a subgroup of \textbf{B}(n), such that
\(\lvert(n, k)\rvert = 2^k\), where elements are binary strings, called \textit{code words}, of 
length n. Code words are the only strings allowed to be transmitted. Due to channel errors where the 
bits flip to the opposite state during transmission, it is not always guranteed that the same code word
is received by the client, so the \textit{received words} are elements of a broader set of all binary 
strings of length n, \textbf{B}(n). 

Consider \(C = \{(0,0,0,0), (0,1,0,1), (1,0,1,0), (1,1,1,1)\}\). For all \(a, b \in C\),
\(a - b \in C\), which establishes that C is a subgroup of \textbf{B}(4). More specifically, \(|C| = 2^2\)
so therefore C is a (4, 2) block code. The received word for C can be any binary string of
length 4, i.e. any element in \textbf{B}(4) = \{0000, 0001, 0010, 0011, 0100, 0101, 0110, 0111, 1000,
1001, 1010, 1011, 1100, 1101, 1110, 1111\}.

%-----------------------------  1 (c & d)  -----------------------------

To systematically generate code words in a block code, we use a \textit{generator matrix}. 
A generator matrix for an (n, k) block code C is a \(k \times n\) matrix G whose rows form 
a basis for C, i.e. are linearly independent. For instance, the generator matrix G, for a (5, 2)
block code \(C = \{(0,0,0,0,0), (1,0,1,1,0), (0,1,1,0,1), (1,1,0,1,1)\}\) is 
\[G = \begin{bmatrix} 
    1 & 0 & 1 & 1 & 0 \\
    0 & 1 & 1 & 0 & 1 \\
\end{bmatrix}\]
To verify if a received word is a valid code word, we make use of a \((n-k) \times n\) 
\textit{parity-check matrix}, \(\mathbf{H}\). The valid code words are the null space of the parity-check 
matrix, i.e., \[HG^T = 0\]
Similarly, a received word C is a valid code word if it satisfies the equation \( HC^T = 0 \).
The parity-check matrix, H, for our (5, 2) block code is

\[
H = \begin{bmatrix}
1 & 0 & 0 & 1 & 0 \\
0 & 1 & 0 & 0 & 1 \\
0 & 0 & 1 & 1 & 1
\end{bmatrix}\]

\[
HG^T = \begin{bmatrix}
1 & 0 & 0 & 1 & 0 \\
0 & 1 & 0 & 0 & 1 \\
0 & 0 & 1 & 1 & 1
\end{bmatrix}
\begin{bmatrix}
1 & 0 \\
0 & 1 \\
1 & 1 \\
1 & 0 \\
0 & 1
\end{bmatrix}
= 
\begin{bmatrix}
0 & 0 \\
0 & 0 \\
0 & 0
\end{bmatrix}
\]

%-----------------------------  1 (e)  -----------------------------
For a given block code, we must define it's capabilities in terms of error detection and error
correction. This capability is defined in terms of the \textit{weight} of a code word in a 
block code. The weight of a binary string \(b \in \mathbf{B}(n)\), denoted by \(w(b)\), is the count of
1s in \(b\). For instance, the weight of \((1,0,1,1,0,1,1,0,1,1,0,1) \in \mathbf{B}(12)\) is 8.
%-----------------------------  1 (f)  -----------------------------
The \textit{distance} D, between two binary strings, \(x , y \in \mathbf{B}(n)\), is the weight of their 
sum in \(\mathbb{Z}_2\), i.e., 
\[D(x, y) = w(x \oplus y)\]
In other words, the distance is the number of positions where the two strings differ.
The distance defines a metric on \textbf{B}(n). Formally, a metric defined on a set is a function such 
that it follows
\begin{itemize}
\item Non-degeneracy: \(d(x, y) = 0 \iff x = y\)
\item Symmetry: \(d(x, y) = d(y, x)\)
\item Triangle inequality: \(d(x, z) \leq d(x, y) + d(y, z)\)
\end{itemize}
It is straightforward to verify that \(D\) satisfies each of these.
%Non-degeneracy holds since 
%\(x \oplus y = 0\) if and only if \(x = y\). Symmetry follows immediately from the commutativity 
%of \(\oplus\). The triangle inequality, which states that flipping bits from \(x\) to \(z\) takes
%no more steps than going via an intermediate string \(y\), is a consequence of the subadditivity
%of weight: \(w(a \oplus b) \leq w(a) + w(b)\) for all \(a, b \in \mathbf{B}(n)\). 
\begin{itemize}
    \item Non-degeneracy: \(D(x,y) = w(x \oplus y) = 0\) if and only if \(x \oplus y = 0\), which holds 
        if and only if \(x_i = y_i\) for all \(i\), i.e. \(x = y\).

    \item Symmetry: \(D(x,y) = w(x \oplus y) = w(y \oplus x) = D(y,x)\), since \(x_i \oplus y_i = y_i 
        \oplus x_i\) for each bit.

    \item Triangle inequality: For any \(x, y, z \in \mathbf{B}(n)\), note that \(x \oplus z = (x \oplus y)
        \oplus (y \oplus z)\), so
        \begin{align*}
            D(x,z) &= w(x\oplus z) \\
                   &= w\bigl((x\oplus y)\oplus(y\oplus z)\bigr) \leq w(x\oplus y) + w(y\oplus z) \\
                   &= D(x,y)+D(y,z) \\
        \end{align*}
        where the inequality uses subadditivity of weight. To see subadditivity holds, observe that
        \((a \oplus b)_i = 1\) only if \(a_i = 1\) or \(b_i = 1\), so every bit counted 
        in \(w(a\oplus b)\) is counted in \(w(a)+w(b)\).
\end{itemize}
This metric structure underlies the error-correcting power of block codes, which we explore further
in the following section.

\section{Decoding and Errors}
%-----------------------------  2 (a)  -----------------------------
Let \(C\) be a code word to be transmitted, whilst \(R\) is its received counterpart. The error in
transmission is measured by the distance between the sent binary string and the received 
string, \(D(C, R)\). These errors change the message seemingly at random coordinates, and therefore,
there is a non zero chance for the received message to have been transformed to another valid code 
word. Let \(d_\text{min}\) be the smallest distance between any two valid code words in a block code,
then the decoder can detect upto \(d_\text{min} - 1\) errors without coinciding with another code word.
If R is detected to have errors, it may be possible to error correct the message. We simply do this by
taking the nearest valid code word to R as the intended message. 

If there is more than one nearest valid code word, the decoder cannot correct the message. 
Therefore, for there to be a unique valid code word for R, R must not be equidistant to two code 
words. Conversely, R must be at most \(\lfloor\frac{d_\text{min} - 1}{2}\rfloor\) from a valid code word.
Therefore, a code with minimum distance \(d_\text{min}\) can correct at most 
\(\lfloor\frac{d_\text{min} - 1}{2}\rfloor\) errors, since any R within that radius is guaranteed to be 
closer to exactly one valid code word than any other. In other words, \(d_\text{min} \geq 2t + 1\), where 
\(t\) is the number of errors the block code can correct.

%-----------------------------  2 (b)  -----------------------------
Let G be a block code described by
\[
G = \begin{bmatrix}
    1 & 0 & 0 & 1 & 1 & 1 \\
    0 & 1 & 0 & 1 & 0 & 1 \\
    0 & 0 & 1 & 1 & 1 & 0 \\
\end{bmatrix}
\]
where the code words for G is defined by the encoding algorithm 
\begin{gather*}
\mathbf{C} = \{B\mathbf{G} : B \in \mathbf{B}(3)\} \subset \mathbf{B}(6) \\
           = \{000000,\ 100111,\ 010101,\ 001110,\ 110010,\ 101001,\ 011011,\ 111100\}
\end{gather*}
Since \(\mathbf{C}\) is closed under addition in \(\mathbb{Z}_2\) and contains 000000, we call C to 
be a linear code. As a result, for any two code words \(C_1, C_2 \in \mathbf{C}\), their sum 
\(C_1 \oplus C_2 \in \mathbf{C}\), and the distance satisfies \(D(C_1, C_2) = w(C_1 \oplus C_2)\).
Therefore, the smallest distance between any two valid code words, \(d_\text{min}\) would be the 
smallest weight of a codeword in C excluding 000000, which by inspection we can deduce to be 3, the 
weight of 010101. Using \(d_\text{min}\), we can infer that \(\mathbf{G}\) can detect at most 
\(d_\text{min} - 1 = 2\) errors and correct \(\lfloor\frac{d_\text{min} - 1}{2}\rfloor = 1\)
error.

%-----------------------------  2 (c) -----------------------------
So far we have defined that the parity-check matrix, \(\mathbf{H}\), detects if an input, \(B \in 
\mathbf{B}(n)\), is valid or not by \(F: \mathbf{B}(n) \rightarrow \mathbf{B}(n - k)\) given by
\[
    F(B) = \mathbf{H}B^t, B \in \mathbf{B}(n)
\]
And if an error is detected, we take the nearest valid code word to be the message. However, to search 
for the nearest valid code word by individually measuring the distance between every possible word 
would be computationally expensive. To solve this issue, we use the fact that \(F\) is a surjective
group homomorphism. We know that F is a group homomorphism, since it satisfies the condition
\(F(a + b) = F(a) + F(b)\), which follows from the linearity of matrix multiplication. We also know that
\(\mathbf{H}\) has a full row rank \(n - k\) which follows from linear independence of \(\mathbf{G}\). 
This suggests that the linear map \(F\) can produce any vector in \(\mathbf{B}(n - k)\) and therefore is
a surjective map. Thus, we can deduce that F is a surjective group homomorphism, which allows us to 
partition \(\mathbf{B}(n)\) into cosets of the kernel of F.
\[
F(B) = HB^t = H(C + e)^t = HC^t + He^t = 0 + He^t = He^t
\]
where e is the error in \(B\). We precompute all cosets of the kernel of \(F\) in \(\mathbf{B}(n)\) and 
find the element with the smallest weight in each coset, called the coset leader, \(\hat{e}\). When 
decoding, we find the coset in which \(F(B) = He^t\) belongs to in our precomputed set of cosets. Once 
we have our coset, and subsequently it's corresponding coset leader \(\hat{e}\), finding the closest 
code word \(\hat{C}\) becomes trivial
\[
    \hat{C} = B - \hat{e}
\]
Since \(F\) is a group homomorphism, its kernel \(ker(F) = \{B \in \mathbf{B}(n):\mathbf{H}B^t = 0\}\) is
a subgroup of \(\mathbf{B}(n)\). Because \(\mathbf{HG}^t\) was defined to be 0, all linear combinations
of \(\mathbf{G}\), i.e. all valid code words, is the kernel of \(F\).

%-----------------------------  2 (d) -----------------------------
A natural question to ask at this point is what is the optimal value for \(k\) such that we maximise
error-control capability whilst minimizing the number of redundant parity bits. This balance is what 
a \textit{Hamming code} achieves. A Hamming code is a linear block code constructed so that the minimum 
distance between any two valid code words is 3, guaranteeing the ability to detect up to 2 errors and 
correct 1 error. This means the coset leaders are precisely the \(n\) single-bit error vectors together 
with the zero vector, giving \(n + 1 = 2^m\) cosets partitioning \(\mathbf{B}(n)\).

For a Hamming code, the parity-check matrix \(H\) is an \(m \times n\) matrix whose columns consist of all 
\(2^m - 1\) distinct non-zero binary vectors of length \(m\). A received word \(B\) is a valid 
codeword if and only if \(HB^T = 0\). If a single-bit error occurs at position \(i\), the syndrome \(HB^t\)
equals precisely the \(i\)-th column of \(H\), which uniquely identifies the coset containing \(B\). Since
each single-bit error vector is the coset leader of a distinct coset, guaranteed by the fact that all 
columns of \(H\) are distinct and non-zero, we can recover \(\hat{e}\) directly from the syndrome, making
the lookup into our precomputed coset table trivial.

\section{Emperical Investigation}
TODO

\end{document}
